\section{Supplementary Material}
\captionsetup{hypcap=false}

This appendix collects details that support the main \benchmark{} claim without
expanding the primary review paper beyond the five-group comparison. The
sections are organized by purpose: reproducibility setup, additional
quantitative checks, and qualitative/scoring details. The appendix is not meant
to introduce a second story; it records the controls behind the same bounded
conclusion in the main paper.

\subsection{Reproducibility Setup and Artifacts}

All trained variants use the same backbone, optimizer objective, adapter
configuration, and distributed setup. The only experimental variable is the
composition of preference rows: answer rows, verifier rows, or their mixture.
This control is important because the paper interprets differences between
groups as evidence about supervision format rather than architecture,
optimizer, or decoding changes.

\begin{center}
\begin{minipage}{0.58\linewidth}
\centering
\footnotesize
\begin{tabular}{ll}
\toprule
Setting & Value \\
\midrule
Backbone & Qwen2.5-VL-7B-Instruct \\
Training objective & Direct Preference Optimization \\
Adapter & LoRA \\
Epochs & 1 \\
DPO beta & 0.1 \\
LoRA rank & 16 \\
Answer rows & 0 or 6{,}000 \\
Verifier rows & 0, 500, 1{,}000, or 2{,}000 \\
Distributed setup & ZeRO-2 \\
\bottomrule
\end{tabular}
\captionof{table}{Shared training setup for the preference-tuned groups. The
main paper reports five groups; the 500-verifier-row run is retained only as
an auxiliary sensitivity check.}
\label{tab:app_train}
\end{minipage}
\end{center}

The completed artifacts are small enough to inspect and regenerate. The
canonical \dataset{} file contains 6{,}000 rows. The selected \dualtwok{} DPO
export contains 8{,}000 rows: 6{,}000 answer-preference rows plus 2{,}000
verifier rows. The evidence probes are fixed at 600 supported-evidence rows
and 400 wrong-evidence rows. This fixed-probe design prevents the ablation from
changing both the training data and the evaluation set at the same time.

Lightweight local regeneration covers data checking, audit-sheet completion,
CEPO-Dual export, and LLaMA-Factory data preparation. Heavy training and full
inference are submitted through Slurm, following the repository rule that GPU
jobs should not be run directly in the interactive environment. The relevant
Slurm entry points are:
\begin{quote}
\footnotesize
\texttt{submit\_cepo\_dual\_pipeline.sh}\\
\texttt{submit\_cepo\_ablation\_pipeline.sh}
\end{quote}
The paper numbers are then refreshed with the deterministic scoring and
summary scripts:
\begin{quote}
\footnotesize
\texttt{score\_object\_eval.py}\\
\texttt{score\_cepo\_evidence\_probe.py}\\
\texttt{bootstrap\_cepo\_probe\_ci.py}\\
\texttt{summarize\_cepo\_ablation\_results.py}
\end{quote}

\subsection{Additional Quantitative Checks}

The main paper keeps the primary comparison to five groups, as required by the
submission plan. Table~\ref{tab:app_ablation} includes the auxiliary
CEPO-Dual-500 run that was useful during analysis but intentionally omitted
from the main body. The table supports two readings. First, verifier-only
supervision improves wrong-evidence rejection but does not replace answer
preference tuning for ordinary short-answer transfer. Second, the mixed
CEPO-Dual variants need enough verifier data before the wrong-evidence probe
clearly improves: CEPO-Dual-500 and CEPO-Dual-1k are weaker than the selected
CEPO-Dual-2k setting.

\begin{center}
\begin{minipage}{\linewidth}
\centering
\scriptsize
\setlength{\tabcolsep}{3.5pt}
\begin{tabular}{@{}lrrrrrrrrr@{}}
\toprule
Setting & Ans. rows & Ver. rows & COCO & GQA & Hard & Supp. & Wrong Rej. & Obj/Attr/Rel & Parse Fail \\
\midrule
Base Instruct & 0 & 0 & 96.0 & 76.2 & 94.4 & 86.5 & 44.3 & 59.5 / 53.7 / 17.6 & 0.0 \\
\answerdpo{} & 6{,}000 & 0 & 97.0 & 76.8 & 95.0 & 90.7 & 34.3 & 45.2 / 45.0 / 10.4 & 0.0 \\
\evidencedpo{} & 0 & 2{,}000 & 95.9 & 76.7 & 94.4 & 87.8 & 53.0 & 77.0 / 63.8 / 16.0 & 0.0 \\
CEPO-Dual-500 & 6{,}000 & 500 & 96.9 & 76.8 & 95.0 & 91.3 & 39.5 & 56.3 / 50.3 / 9.6 & 0.0 \\
\dualonek{} & 6{,}000 & 1{,}000 & 97.0 & 77.2 & 94.9 & 93.2 & 48.5 & 74.6 / 59.1 / 9.6 & 0.0 \\
\dualtwok{} & 6{,}000 & 2{,}000 & 97.0 & 77.1 & 94.9 & 95.2 & 70.3 & 96.0 / 85.9 / 25.6 & 0.0 \\
\bottomrule
\end{tabular}
\captionof{table}{Full verifier-count ablation from the released artifacts.
The main paper omits CEPO-Dual-500 to keep the primary comparison to at most
five groups.}
\label{tab:app_ablation}
\end{minipage}
\end{center}

The ablation also shows why the paper does not present CEPO-Dual as a broad
hallucination method. The largest gain is on the internal evidence probe, not
on every external yes/no benchmark. Table~\ref{tab:app_external} therefore
serves as a sanity check. CEPO-Dual-2k stays close to CEPO Answer-DPO on POPE
and AMBER, which means the verifier rows do not cause an obvious collapse in
ordinary polling behavior. The differences are small, so the external table is
supporting evidence for stability rather than the source of the main claim.

\begin{center}
\begin{minipage}{0.55\linewidth}
\centering
\footnotesize
\begin{tabular}{lccc}
\toprule
Eval & Base & \answerdpo{} & \dualtwok{} \\
\midrule
POPE random & 88.7 & 89.6 & \textbf{89.8} \\
POPE popular & 87.7 & 88.5 & \textbf{88.7} \\
POPE adversarial & 86.7 & 87.1 & \textbf{87.3} \\
AMBER discr. & 87.6 & \textbf{88.2} & \textbf{88.2} \\
\bottomrule
\end{tabular}
\captionof{table}{External yes/no sanity checks. Values are accuracy
percentages under the same deterministic parser.}
\label{tab:app_external}
\end{minipage}
\end{center}

\subsection{Qualitative Analysis and Scoring Rules}

The most persistent failure mode is relation verification. Object and
attribute evidence can often be checked by recognizing one label or property,
but a left/right relation requires the model to preserve subject, object, and
direction at the same time. Table~\ref{tab:app_qual} gives representative
probe outputs. The successful wrong-relation case shows that CEPO-Dual-2k can
learn to reject swapped evidence, while the failure case shows the remaining
weakness: if both objects are visible, the model may still accept the swapped
relation as supported.

\begin{center}
\begin{minipage}{\linewidth}
\centering
\scriptsize
\begin{tabular}{@{}p{0.14\linewidth}p{0.34\linewidth}p{0.16\linewidth}p{0.13\linewidth}p{0.10\linewidth}@{}}
\toprule
Type & Candidate claim and evidence & Target & \dualtwok{} output & Reading \\
\midrule
Supported attr. &
claim: vehicle is white; evidence: label=vehicle, attribute=white &
supported & supported & correct \\
Supported material &
claim: window is glass; evidence: label=window, attribute=glass &
supported & supported & correct \\
Wrong object &
claim: sink visible; evidence: label=potted plant &
wrong evidence & wrong evidence & correct \\
Wrong relation &
claim: glass left\_of grape; evidence swaps subject/object &
wrong evidence & contradicted & correct rejection \\
Wrong relation &
claim: man right\_of parking meter; evidence swaps subject/object &
wrong evidence & supported & relation failure \\
\bottomrule
\end{tabular}
\captionof{table}{Representative probe outputs. Examples are anonymized to
local image references and show why relation reversals remain the hardest
slice.}
\label{tab:app_qual}
\end{minipage}
\end{center}

Scoring follows the same conservative principle. Short-answer metrics use a
deterministic yes/no parser. Outputs with explicit uncertainty such as
``cannot determine'' or ``not sure'' are refusals; outputs without a parseable
yes/no decision are counted as other/invalid. The evidence-probe scorer first
tries to parse JSON and then falls back to labeled fields. For wrong-evidence
rows, both \texttt{wrong\_evidence} and \texttt{contradicted} count as
rejecting the candidate evidence, because both labels indicate that the
evidence should not be accepted as support for the claim. This scoring rule is
why the paper reports parser quality next to the main probe metrics: an
evidence-verification gain should not come from malformed output or a fragile
parser recovery.
